\textit{Le service forestier a essayé de faciliter leur vie dans ces petites communautés où la
vie n’était pas toujours drôle..}

\textit{Finalement c’est en 1965, après 3 ans de séjour à Meyrueis, qu’ils furent relogés
dans des appartements plus modernes et mieux équipés à Mende (Sonacotra) et à
Lodève.
Il y eut même un projet, beaucoup plus tard, en 1980, de tous les regrouper en
Provence, puis de les transférer - une idée de technocrates parisiens du Ministère
de l’Agriculture - en Haute Corse dans une des zones les plus isolées et difficile
d’accès. Il n’y eut qu’un cri : pas question. Le service forestier comme le préfet
prirent position en leur faveur. Et rien ne se passa. Malgré toutes les vicissitudes
qu’ils connurent, leur état d’esprit m’apparaissait assez serein.}

\textit{Et maintenant ? 25 ans plus tard ?}

\textit{A Mende plusieurs familles se sont installées définitivement et ont fait souche. Les
rapports avec la population locale n’ont pas posé de gros problèmes ; il y a un
Conseiller municipal, des sportifs, (foot, coureurs ...) Au début le taux de chômage
était élevé mais petit à petit des employeurs locaux en ont embauché, la majorité
d’entre eux se sont casés ou ont même lancé de petites affaires.
Les harkis d’origine étaient assimilés à des employés de l’ONF avec droit à la
retraite. Ils ont demandé que leurs fils disposent du même statut mais ce leur fut
refusé par crainte de « ghettoisation. »}

\begin{center}
\textbf{En guise de conclusion provisoire}
\end{center}

Aux dires des témoins, il semble que, même s’il n’y eut pas beaucoup de contacts
avec la population de Meyrueis - mais on sortait tout juste de la guerre d’Algérie- les
choses se passèrent moins mal, dans ces petits « hameaux », que dans les grands
camps comme Bias en Lot et Garonne, Saint Maurice l’Ardoise dans le Gard ou
Rivesaltes dans l’Aude, où furent hébergées des centaines de familles de harkis avec
une discipline de fer, des barbelés et des portails fermés la nuit, des prénoms français
donnés d’autorité par les chefs de camp aux bébés venant de naître etc...
Malheureusement je n’ai pu, à ce jour, retrouver d’anciens harkis de la Foulatière ni
leurs femmes ni leurs enfants...

\begin{flushright}
\textit{A Meyrueis\\
Le 26 Novembre 2006}
\end{flushright}

Avant de rassembler et photocopier ces documents le 10 Décembre 2006, j’ai pu jeter
un coup d’œil sur des rapports de responsables du hameau de Meyrueis. J’en ignorais
l’existence jusqu’alors ! j’ai pu, également, avoir les coordonnées d’un fils de harki
de Villemagne, installé à Mende, mais pas encore celles de familles ayant séjourné à
Meyrueis. Voici cependant quelques renseignements complémentaires : notés dans ces
rapports ou d’après la très brève conversation téléphonique avec M. S, fils de harki.
